% ============================================================
% 04-experiments.tex
% ============================================================
\section{Experimental Results}
\label{sec:exp}

This section presents the empirical evaluation of \system{}.

%------------------------------------------------
\subsection{Experimental Benchmark and Protocols}
\label{sec:dataset}

To support empirical evaluation across representative counseling scenarios, our evaluation uses a corpus of $200$ administrative queries constructed via a document-grounded semi-automated protocol over official Can Tho University (CTU) statutes. Candidate queries covering diverse conversational phrasing (e.g., student colloquialisms, prerequisite rules, and fee schedules) were generated from statutory articles and reviewed against the institutional corpus to verify ground-truth intents, gold passages, and factual answers. The corpus is partitioned into a \textbf{Development Benchmark ($N=100$ queries)}, used for offline index calibration, prompt engineering, and qualitative failure analysis, and a separately constructed \textbf{Held-Out Test Set ($N=100$ queries)}. The held-out set follows a predefined complexity distribution comprising 40 direct single-hop, 20 within-domain multi-hop, 20 cross-domain, 10 entity-comparison, 5 temporal, and 5 adversarial queries. It contains 51 formal and 49 colloquial formulations; 56 queries have one gold source and 44 require multiple gold sources.

The construction audit used entity inventories separated from those in the development benchmark. It found no exact duplicate queries across the two splits and recorded a maximum cross-split 5-gram Jaccard similarity of 0.1875. Gold sources and required facts were reviewed against 244 institutional documents and 559 structured tuition records. The four single-domain subsets contain 19 Academic, 28 Financial, 17 Scholarship, and 16 General queries; the remaining 20 queries combine evidence across domains and are therefore analyzed separately rather than assigned to a single specialist domain. All primary metrics in Tables~\ref{tab:retrieval_evaluation} and~\ref{tab:e2e_evaluation} are evaluated on this held-out set.

\noindent\textbf{Execution and Evaluation Invariants:} Generative and automated evaluation tasks were executed using Google Gemini 2.5 Flash Lite via Google Cloud Vertex AI under zero-temperature greedy decoding ($T=0.0$). To account for subtle API-level non-determinism across cloud endpoints, each query was evaluated across $R=3$ repeated evaluation runs, yielding $2,100$ end-to-end generation and Ragas evaluation turns ($7\text{ configurations} \times 100\text{ queries} \times 3\text{ runs}$).

\noindent\textbf{Computational Infrastructure:} Neural cross-encoder reranking was executed on an NVIDIA RTX 1060 GPU (4~GB VRAM). Relational curricula are hosted in Neo4j 5 Enterprise with APOC extensions; document parent chunks (2,800 characters, 100 overlap) reside in PostgreSQL, while child chunks (400 characters, 50 overlap) are indexed in Qdrant with \texttt{vietnamese-bi-encoder} (768-d). Neural candidate reranking uses \texttt{bge-reranker-v2-m3}. Multi-agent coordination was orchestrated via LangGraph under a StateGraph supervisor with Redis 7 caching.

%------------------------------------------------
\subsection{Evaluation Metrics}
\label{sec:metrics}

\noindent\textbf{Retrieval and Multi-Agent Benchmarks:} Retrieval effectiveness (Scenario~1) is measured via standard ranking metrics on top-$k$ returned source documents: Hit@$k$ ($k \in \{1, 3\}$), Precision@5, Recall@5, and Mean Reciprocal Rank (MRR@10); domain-stratified analysis reports Hit@1. Multi-agent coordination (Scenario~3) is evaluated across isolated operational decisions: Routing Agent Accuracy and Macro-F1 across the four domains, Intent Classification Accuracy, Tool Selection and Argument Exact Match (EM), End-to-End Pass Rate, Adversarial Tool-Gating Robustness, and Safe Result Behavior under strict operational limits ($\text{max\_tool\_calls}=1, \text{timeout}=60\,\text{s}$).

\noindent\textbf{End-to-End Generative Evaluation:} End-to-end synthesis (Scenario~2) is evaluated using the Ragas framework~\cite{es2023ragas} alongside ground-truth regulatory alignment over Top-7 context: Answer Relevancy (AR), Context Recall (CR), Context Precision (CP), and Answer Correctness (AC). In addition, we evaluate Factual Exact Match (Fact EM) for strict numerical tuition fee and credit agreement, Source Recall for gold-source coverage, and Source Average Precision (Source AP) for verified regulatory citations. Evaluating generative QA through LLM-as-a-judge is well-established~\cite{zheng2023judging,es2023ragas}; to mitigate evaluation bias, our framework enforces reference-grounded metrics conditioned on gold documents, greedy decoding ($T\!=\!0.0$) across $R\!=\!3$ repeated runs, and a programmatic Fact EM evaluator for tuition arithmetic.

%------------------------------------------------
\subsection{Experimental Scenarios and Benchmark Results}
\label{sec:scenarios}

To address the three research questions, we evaluate \system{} across three coordinated experimental scenarios: retrieval component stacking (RQ3), end-to-end QA generation and modular ablation (RQ2), and multi-agent coordination and tool reliability (RQ1).

\subsubsection{Scenario 1 --- Retrieval Component Stacking (E1--E5):}
\emph{Objective \& RQ Alignment:} Scenario~1 directly addresses \textbf{RQ3} by comparing five retrieval configurations at aggregate and complexity-stratified levels on the 100-query Held-Out Test Set.
\emph{Configurations:} We compare five cumulative configurations (E1--E5, detailed in Table~\ref{tab:retrieval_evaluation}) transitioning from single sparse/dense baselines to hybrid fusion, neural cross-reranking, and our proposed intent-guided evidence stack.

\emph{Results:} As shown in Table~\ref{tab:retrieval_evaluation}, E5 obtains the highest values across the five reported retrieval metrics, reaching 0.7500 Hit@1, 0.9100 Hit@3, 0.2380 Precision@5, 0.8458 Recall@5, and 0.8354 MRR@10. Relative to E4, E5 increases Hit@1 by 7.0 percentage points and MRR@10 by 0.0450; relative to E1, the Hit@1 difference is 22.0 percentage points. These improvements are accompanied by an increase in mean retrieval latency from 15.62~s for E4 to 17.09~s for E5. Because E5 jointly adds intent governance, graph grounding, and parent--child mapping, the E4-to-E5 contrast represents the integrated bundle rather than an isolated routing or graph effect.

\begin{table}[!htbp]
\centering
\caption{Scenario 1: Progressive retrieval component stacking on the redesigned Held-Out Test Set ($N=100$).}
\label{tab:retrieval_evaluation}
\resizebox{\columnwidth}{!}{
\begin{tabular}{lcccccc}
\toprule
\textbf{Retrieval Configuration} & \textbf{Hit@1} & \textbf{Hit@3} & \textbf{P@5} & \textbf{Recall@5} & \textbf{MRR@10} & \textbf{Latency (ms)} \\
\midrule
E1: BM25 Lexical Baseline        & 0.5300 & 0.7600 & 0.1980 & 0.7133 & 0.6661 & \textbf{9.79} \\
E2: Dense Vector Semantic        & 0.4200 & 0.6300 & 0.1540 & 0.5542 & 0.5223 & 50.80 \\
E3: Hybrid RRF (BM25 + Dense)    & 0.5000 & 0.7800 & 0.2040 & 0.7308 & 0.6568 & 53.10 \\
E4: Hybrid + Neural Reranker     & 0.6800 & 0.9000 & 0.2320 & 0.8275 & 0.7904 & 15,624.33 \\
\textbf{E5: CTU-Chat Proposed Retrieval} & \textbf{0.7500} & \textbf{0.9100} & \textbf{0.2380} & \textbf{0.8458} & \textbf{0.8354} & 17,094.94 \\
\bottomrule
\end{tabular}
}
\end{table}

\begin{table}[!htbp]
\centering
\caption{E4--E5 retrieval results by held-out query complexity.}
\label{tab:retrieval_complexity}
\begingroup
\setlength{\tabcolsep}{7pt}
\renewcommand{\arraystretch}{1.12}
\begin{tabular}{lccccc}
\toprule
\multirow{2}{*}{\textbf{Complexity Tier}} & \multirow{2}{*}{\textbf{$N$}} & \multicolumn{2}{c}{\textbf{Hit@1}} & \multicolumn{2}{c}{\textbf{Recall@5}} \\
\cmidrule(lr){3-4}\cmidrule(lr){5-6}
& & \textbf{E4} & \textbf{E5} & \textbf{E4} & \textbf{E5} \\
\midrule
Direct single-hop       & 40 & 0.6500 & \textbf{0.6750} & 0.9000 & \textbf{0.9250} \\
Within-domain multi-hop & 20 & 0.7000 & 0.7000 & 0.7167 & \textbf{0.7417} \\
Cross-domain            & 20 & 0.7000 & \textbf{0.8500} & \textbf{0.7833} & 0.7750 \\
Entity comparison       & 10 & 0.8000 & \textbf{1.0000} & 0.9500 & 0.9500 \\
Temporal                 & 5 & 0.8000 & \textbf{1.0000} & 1.0000 & 1.0000 \\
Adversarial              & 5 & 0.4000 & 0.4000 & 0.4500 & \textbf{0.5500} \\
\bottomrule
\end{tabular}
\endgroup
\end{table}

The E4-to-E5 Hit@1 differences are largest on entity-comparison and temporal queries (20 percentage points each) and cross-domain queries (15 percentage points). The two configurations tie on multi-hop and adversarial Hit@1, while E5 obtains higher Recall@5 on both tiers. On cross-domain queries, however, E5 has slightly lower Recall@5 than E4 (0.7750 versus 0.7833), indicating that its improvement in first-rank placement does not extend uniformly to evidence coverage.

\begin{table}[!htbp]
\centering
\caption{Scenario 2: End-to-end QA generation and modular ablation on the redesigned Held-Out Test Set ($N=100, R=3$, 2,100 total evaluation turns, Top-7 context).}
\label{tab:e2e_evaluation}
\small
\setlength{\tabcolsep}{3.5pt}
\begin{tabular}{lccccccc}
\toprule
\textbf{Cfg.} & \textbf{CP} & \textbf{CR} & \textbf{AR} & \textbf{AC} & \textbf{Fact EM} & \textbf{Src. R} & \textbf{Src. AP} \\
\midrule
T1 & 0.2576 & 0.5411 & 0.3961 & 0.4758 & 0.5816 & 0.7308 & 0.5498 \\
T2 & 0.1664 & 0.2682 & 0.2281 & 0.3607 & 0.5443 & 0.5642 & 0.4240 \\
T3 & 0.2110 & 0.4843 & 0.3618 & 0.4631 & 0.5965 & 0.7608 & 0.5469 \\
\textbf{T4} & 0.3833 & 0.6296 & 0.4454 & 0.5080 & \textbf{0.6529} & \textbf{0.8408} & \textbf{0.7313} \\
T5 & 0.2255 & 0.4617 & 0.2996 & 0.4400 & 0.5857 & 0.7192 & 0.6104 \\
T6 & 0.3761 & 0.5979 & \textbf{0.4639} & \textbf{0.5108} & 0.6271 & \textbf{0.8408} & 0.7058 \\
T7 & \textbf{0.3862} & \textbf{0.6572} & 0.4235 & 0.4992 & 0.6384 & 0.8208 & 0.7238 \\
\bottomrule
\end{tabular}
\smallskip
\parbox{\columnwidth}{\footnotesize\emph{Configurations:} T1: BM25; T2: dense; T3: hybrid RRF; T4: full system; T5: without cross-reranker; T6: without Neo4j grounding; T7: without subspace governance. \emph{Metrics:} Fact EM: Factual Exact Match on tuition/credits; CP: Context Precision; CR: Context Recall; AR: Answer Relevancy; AC: Answer Correctness; Src. R: Source Recall; Src. AP: Source Average Precision. Ragas values are means over 300 query--repetition observations per configuration; corresponding standard deviations and bootstrap confidence intervals are retained in the experiment artifacts. All metrics are normalized to $[0,1]$; higher is better.}
\end{table}

%------------------------------------------------
\subsubsection{Scenario 2 --- End-to-End QA Generation and Modular Ablation (T1--T7):}
\emph{Objective \& RQ Alignment:} Scenario~2 directly addresses \textbf{RQ2} by benchmarking the proposed full system (T4) against three single-agent RAG baselines (T1--T3) and three modular ablations (T5--T7) on the Held-Out Test Set ($2,100$ total evaluation turns; Table~\ref{tab:e2e_evaluation}).
\emph{Results \& Interpretation:} As reported in Table~\ref{tab:e2e_evaluation}, T4 obtains higher values than the BM25 baseline across all seven reported metrics, including a 7.13-percentage-point difference in Fact~EM ($65.29\%$ vs.\ $58.16\%$). Across all configurations, T4 achieves the highest Fact~EM and Source AP and ties T6 for the highest Source Recall. T6 records the highest Answer Relevancy and Answer Correctness, while T7 records the highest Context Precision and Context Recall. The ablation results therefore indicate metric-specific differences rather than uniform dominance by a single configuration.
\FloatBarrier

%------------------------------------------------
\subsubsection{Scenario 3 --- Multi-Agent Routing, Tool Reliability, and Robustness:}
\emph{Objective \& RQ Alignment:} Scenario~3 directly addresses \textbf{RQ1} by evaluating multi-agent coordination, tool execution reliability, and bounded safety independently from retrieval or generation.
\emph{Experimental Setup:} As reported in Table~\ref{tab:scenario3_results}, evaluation spans three suites ($N=550$ runs): (1)~\textbf{Supervisor Routing (Panel A):} 100 single-domain benchmark queries ($N=300$ for LLM Supervisor vs. $N=100$ for Rule-Based Router); (2)~\textbf{Specialist Tool Reliability (Panel B):} 30 test cases ($N=90$ runs across 3 repeated runs) covering structured tuition lookup, scholarship calculation, and deterministic tuition reduction arithmetic; and (3)~\textbf{Adversarial Robustness (Panel C):} 20 failure stress tests ($N=60$ runs across 3 repeated runs) challenging boundary invariants.

\begin{table}[!htbp]
\centering
\caption{Scenario 3: Multi-agent routing, specialist tool execution reliability, and robustness under a maximum of one tool call per decision ($\text{max\_tool\_calls}=1$).}
\label{tab:scenario3_results}
\resizebox{\columnwidth}{!}{
\begin{tabular}{lccccccc}
\toprule
\multicolumn{8}{l}{\textbf{Panel A: Supervisor Routing and Specialist Selection ($N=100$ queries; Rule: 1 rep / LLM: 3 reps, $N=300$)}} \\
\midrule
\textbf{Router Variant} & \textbf{Runs ($N$)} & \textbf{Agent Acc.} & \textbf{Macro-F1} & \textbf{Intent Acc.} & \textbf{Bounded} & \multicolumn{2}{c}{\textbf{Routing Paradigm}} \\
\midrule
Rule-Based Router          & 100 & 85.0\% & 69.2\% & 69.0\% & 100.0\% & \multicolumn{2}{c}{Heuristic Regex} \\
LLM Supervisor (\system{}) & 300 & \textbf{97.0\%} & \textbf{97.8\%} & \textbf{95.0\%} & \textbf{100.0\%} & \multicolumn{2}{c}{Structured Pydantic} \\
\midrule
\multicolumn{8}{l}{\textbf{Panel B: Specialist Tool Execution Reliability ($N=90$ runs across 30 test cases, 3 repeated runs)}} \\
\midrule
\textbf{Specialist Tool Function} & \textbf{Runs ($N$)} & \textbf{Selection} & \textbf{Arg. EM} & \textbf{Result Acc.} & \textbf{E2E Pass} & \textbf{Bounded} & \textbf{Paradigm} \\
\midrule
\texttt{tuition\_lookup}                & 30 & 100.0\% & 100.0\% & 100.0\% & 100.0\% & 100.0\% & Structured Catalog \\
\texttt{scholarship\_calculation}       & 30 & 86.7\%  & 76.7\%  & 86.7\%  & 76.7\%  & 100.0\% & LLM Tool Selection \\
\texttt{tuition\_reduction\_calculation}& 30 & 100.0\% & 100.0\% & 100.0\% & 100.0\% & 100.0\% & Arithmetic Tool \\
\midrule
\multicolumn{8}{l}{\textbf{Panel C: Robustness and Bounded Failure Handling ($N=60$ runs under malformed/adversarial inputs)}} \\
\midrule
\textbf{Adversarial Gate} & \textbf{Runs ($N$)} & \textbf{Agent Acc.} & \textbf{Tool Gate} & \textbf{Safe Result} & \textbf{E2E Pass} & \textbf{Bounded} & \textbf{Notes} \\
\midrule
Failure Gate Robustness    & 60 & 100.0\% & 100.0\% & 98.3\% & 98.3\% & 100.0\% & -- \\
\bottomrule
\end{tabular}
}
\smallskip
\parbox{\columnwidth}{\footnotesize\emph{Note:} \texttt{tuition\_lookup} evaluates deterministic structured catalog queries; scholarship and tuition-reduction cases evaluate LLM-driven parameter extraction and external deterministic tool execution.}
\end{table}

\subsection{Discussion}
\label{sec:discussion}

\textbf{Answering RQ1 (Domain-Specialized Multi-Agent Architecture):} Centralizing routing under a structured supervisor with asymmetric privileges and bounds ($\text{max\_tool\_calls}=1$, timeout $=60\,\text{s}$) demonstrated reliable coordination under the evaluated configuration: the supervisor achieves $97.0\%$ agent accuracy, $97.8\%$ Macro-F1, and $95.0\%$ intent accuracy across 300 runs, outperforming rule-based routing ($85.0\%$, $69.2\%$). In specialist tool execution (Panel~B), \texttt{tuition\_lookup} and \texttt{tuition\_reduction\_calculation} achieved $100.0\%$ end-to-end pass rates, while \texttt{scholarship\_calculation} reached $76.7\%$ ($86.7\%$ tool selection), with all three functions maintaining $100.0\%$ bounded completion. Furthermore, adversarial testing (Panel~C) verified $100.0\%$ tool-gating accuracy, $98.3\%$ safe-result behavior, and $100.0\%$ bounded completion. To assess architectural scalability, a paired stress-test compared \system{} against a monolithic single-agent baseline: under the baseline registry (11 tools), both architectures exhibited comparable performance (single-agent $93.3\%$ vs. multi-agent $88.3\%$ E2E, difference not statistically significant); however, when 16 synthetic distractor tools were introduced to simulate institutional expansion (27 tools total), the monolithic agent suffered from decision-space pollution, degrading from $98.3\%$ to $91.7\%$ in tool selection. In contrast, domain-partitioned multi-agent maintained $100.0\%$ tool selection (paired $\Delta = +8.3\%$, 95\% bootstrap CI $[+1.7, +15.0]$, $n=60$). The benefit of domain-specialized multi-agent decomposition became more apparent as the tool registry expanded, because partitioned tool access reduced decision-space interference relative to the evaluated monolithic single-agent baseline.

\textbf{Answering RQ2 (Heterogeneous Knowledge Allocation \& Modular Ablations):} The modular ablations reveal metric-specific effects. Relative to T5, T4 obtains higher scores on every reported metric, including Answer Relevancy ($0.4454$ vs.\ $0.2996$), Context Recall ($0.6296$ vs.\ $0.4617$), Context Precision ($0.3833$ vs.\ $0.2255$), and Fact~EM ($65.29\%$ vs.\ $58.57\%$), supporting the contribution of cross-reranking within the evaluated configuration. The graph comparison is mixed: relative to T6, T4 has higher Context Recall ($0.6296$ vs.\ $0.5979$), Source AP ($0.7313$ vs.\ $0.7058$), and Fact~EM ($65.29\%$ vs.\ $62.71\%$), while T6 has slightly higher Answer Relevancy and Answer Correctness and ties T4 on Source Recall. Similarly, T7 has slightly higher Context Precision and Context Recall than T4, whereas T4 has higher Answer Relevancy, Answer Correctness, Source Recall, Source AP, and Fact~EM. These comparisons support component-specific associations within the evaluated setting but do not establish uniform causal effects across query types or metrics.

\textbf{Answering RQ3 (Representation-Aware Retrieval Effectiveness):} On the redesigned benchmark, E3-to-E4 improvements appear in both early-rank ordering and evidence coverage, while E5 provides smaller additional gains across the aggregate metrics. Complexity stratification shows that E5's largest Hit@1 differences relative to E4 occur on entity-comparison, temporal, and cross-domain queries. Direct and multi-hop queries show smaller or no Hit@1 differences, and the adversarial tier remains difficult for both configurations. Moreover, E5's cross-domain Recall@5 is slightly below E4 despite its higher Hit@1. Because E5 changes intent governance, graph grounding, and parent--child mapping together, these results compare complete configurations and do not attribute the observed gains causally to any one component.

%------------------------------------------------
\subsection{Threats to Validity and Limitations}
\label{sec:limitations}

The evaluation is scoped to Can Tho University; multi-institutional validation remains planned. Although the redesigned held-out set explicitly covers multiple complexity tiers, its temporal and adversarial strata contain only five queries each. In addition, the four-domain analysis covers the 80 single-domain queries and excludes the 20 cross-domain queries, so it should not be interpreted as a complete partition of the benchmark. Gemini 2.5 Flash Lite served as both generator and judge, introducing potential evaluation bias despite reference grounding, deterministic decoding, and repeated runs. Curricula graphs rely on offline schemas, and automated judges may penalize valid abstentions. Fault injection across external infrastructure (Neo4j, Qdrant) is reserved for future work.
